\chapter{Relation with other axiomatic systems}
\label{chap:modcatcatcofobj}

We would like to describe in this chapter how ABC cofibration categories relate to Brown's categories of cofibrant objects, to Quillen model categories and to other axiomatizations that have been proposed for categories with cofibrations. 

Aside from the goal of bringing together and comparing various axiomatizations that have been proposed for (co)fibrations and weak equivalences, this allows us to tap into a large class of examples of ABC model categories. 

For example, simplicial sets $sSets$ form a Quillen model category \cite{Quillen1}, with inclusions as cofibrations, with maps satisfying the Kan extension property as fibrations and with maps whose geometric realization is a topological homotopy equivalence as weak equivalences. Any Quillen model category is an ABC model category, and therefore $sSets$ is an ABC model category. By \theoremref{thm:generalpointwisecofstructure}, so is any diagram category $sSets^{\cat D}$ for a small category ${\cat D}$, and by \definitionref{defn::restricteddiagrams} and \theoremref{thm:generalpointwisecofibstruct2} so are the ${\cat D}_2$-reduced ${\cat D}_1$-diagrams $sSets^{({\cat D}_1, {\cat D}_2)}$ for a small category pair $({\cat D}_1, {\cat D}_2)$.

Other basic examples of ABC model categories are explained in \chapterref{chap:examples}.

The list of alternative cofibration category axiomatizations discussed in this chapter is by no means exhaustive. 

We have omitted Alex Heller's notion of h-c categories \cite{HellerHC1}, \cite{HellerHC2}, \cite{HellerHC3}. As we have seen, the ABC cofibration category axioms are written in terms of {\it cofibrations} and {\it weak equivalences}.  In contrast, Heller's h-c category axioms are written in terms of cofibrations and a homotopy relation $\simeq$ on maps. A map $f \colon A \ra B$ in a h-c category by definition is a weak equivalence if there exists $g \colon B \ra A$ with $gf \simeq 1_A$, $fg \simeq 1_B$.

We have also omitted the categories with a natural cylinder \cite{Kamps}, \cite{Shitanda} and \cite{Kamps-Porter}, which are a surprising elaboration of ideas surrounding the Kan extension property for cubical sets \cite{Kan1}, \cite{Kan2}.



%We would like to mention without proof that the classic topological category $Top$ also forms an ABC model category, with homotopy equivalences as weak equivalences, with cofibrations as being the maps satisfying the {\it homotopy extension property} (HEP) and fibrations satisfying the {\it homotopy lifting  property} (HLP). \footnote{In fact, there is a very intriguing axiomatic approach to the proof of $Top$ being an ABC cofibration category, for which we refer the reader to Kamps \cite{Kamps}, Shitanda \cite{Shitanda} and Kamps-Porter \cite{Kamps-Porter}. In outline, suppose that a category ${\cat M}$ has a natural cylinder functor, defined as a functor $I \colon {\cat M} \ra {\cat M}$ with natural maps $i_0, i_1 \colon A \ra IA$ and $p \colon IA \ra A$, satisfying $pi_0 = pi_1 = 1_A$. The cylinder allows us to define weak equivalences as homotopy equivalences and cofibrations by the homotopy extension property, as we do in $Top$. Then if ${\cat M}$ is cocomplete, the functor $I$ preserves colimits and the {\it cubical} function complexes $Hom_{\cat M}(I \tdot A, B)$ satisfy a number of cubical Kan extension conditions in a way made explicit in \cite{Kamps-Porter}, then ${\cat M}$ is an ABC cofibration category.}

\section{Brown's categories of cofibrant objects}
In his paper \cite{Brown}, Brown defines categories of fibrant objects (and dually categories of cofibrant objects). We list below Brown's axioms, stated in the cofibration setting, slightly modified but equivalent to the actual axioms of \cite{Brown}.

\begin{defn}[Categories of cofibrant objects]
\label{defn:catcofobj}
\index{category!of (co)fibrant objects}
A category of cofibrant objects $({\cat M}$, ${\cat W}$, ${\cat Cof})$ consists of a category ${\cat M}$ and two distinguished classes of maps ${\cat W}, {\cat Cof}$ - the weak equivalences and respectively cofibrations of ${\cat M}$, subject to the axioms below:
\begin{description}
\item[CFObj1] 
All isomophisms of ${\cat M}$ are trivial cofibrations. ${\cat M}$ has an initial object $\initial$, and all objects of ${\cat M}$ are cofibrant. Cofibrations are stable under composition. 
\item[CFObj2] 
(Two out of three axiom) If $f, g$ are maps of ${\cat M}$ such that $gf$ is defined, and if two of $f, g, gf$ are weak equivalences, then so is the third.
\item[CFObj3] 
(Pushout axiom) Given a solid diagram in ${\cat M}$, with $i$ a cofibration,
\begin{center}
$ \xymatrix {
	A \ar[r] \ar@{>->}[d]_i &
	C \ar@{>->}[d]^j \\
	B \ar@{-->}[r] &
	D
	}$
\end{center}
then the pushout exists in ${\cat M}$ and $j$ is a cofibration. If additionally $i$ is a trivial cofibration, then $j$ is a trivial cofibration.
\item[CFObj4]
(Cylinder axiom) For any object $A$ of ${\cat M}$, the codiagonal $\nabla \colon A \Sum A \ra A$ admits a factorization as a cofibration followed by a weak equivalence.
\end{description}
\end{defn}

The axioms for categories of fibrant objects are dual to those of \definitionref{defn:catcofobj}, and are denoted FObj1-FObj4.

The precofibration categories (satisfying the minimal axioms CF1-CF4 but not the additional axioms CF5-CF6) are essentialy a modification of Brown's categories of cofibrant objects - in the sense that we allow objects to be non-cofibrant. The following lemma explains the precise relationship between precofibration categories and categories of cofibrant objects.

\begin{prop}
\label{prop:relcatcofobj}
\mbox{}
\begin{enumerate}
\item
Any category of cofibrant objects is a precofibration category. Conversely, if ${\cat M}$ is a precofibration category, then ${\cat M}_{cof}$ is a category of cofibrant objects. 
\item
Any category of fibrant objects is a prefibration category. Conversely, if ${\cat M}$ is a prefibration category, then ${\cat M}_{fib}$ is a category of fibrant objects. 
\end{enumerate}
\end{prop}

\begin{proof}
We only prove (1). Implication $\Leftarrow$ is an easy consequence of the axioms. For the other direction $\Rightarrow$, the only axiom that needs to be proved is the factorization axiom CF4.

For that, it suffices to prove Brown's factorization lemma \ref{lem:brownfact} in the context of the axioms CFObj1-CFObj4. We want to show that any map $f \colon A \ra B$ factors as $f = rf'$, where $f'$ is a cofibration and $rs = 1_B$ for a trivial cofibration $s$.

Choose a cylinder $IA$, and construct $s$ as the pushout of the trivial cofibration $i_0$.

\begin{center}
$ \xymatrix {
        A \ar[r]^f \ar@{>->}[d]_{i_0}^{\sim} &
	B \ar@{>->}[d]^s_{\sim} & \\
	IA \ar[r]^F &
	B'
	}$
\end{center}

Notice that $i_0$ has $p \colon IA \lra A$ as a left inverse, and it follows that $s$ has a left inverse $r$. Let $f'$ be $Fi_1$, which satisfies $f = rf'$, and to complete the proof it remains to prove that $f'$ is a cofibration. 

We notice that $B'$ is also the pushout of the diagram below
\begin{center}
$ \xymatrix {
        A \Sum A \ar[r]^{f \Sum 1} \ar@{>->}[d]_{i_0 + i_1} &
	B \Sum A \ar@{>->}[d]^{s + f'} & \\
	IA \ar[r]^F &
	B'
	}$
\end{center}
so $f'$ is $ \xymatrix { A \ar@{>->}[r] & B \Sum A \ar@{>->}[r]^-{s + f'} & B' } $, therefore a cofibration.
\end{proof}

\section{Quillen model categories}
\index{Quillen!model category}
Quillen's model categories involve both cofibrations and fibrations, and come with built-in Eckman-Hilton duality between cofibrations and fibrations. 

To start, recall the definition of the left (and right) lifting property of maps.

\begin{defn}
\label{defn:llp}
\index{left lifting property LLP}
\index{right lifting property RLP}
For a solid commutative diagram in a category ${\cat M}$
\begin{center}
$ \xymatrix {
	A \ar[r] \ar@{->}[d]_i &
	C \ar@{->}[d]^p \\
	B \ar[r] \ar@{-->}[ur] &
	D
	}$
\end{center}
if a dotted arrow exists making the diagram commutative we say that $i$ has the LLP (left lifting property) with respect to $p$, and that $p$ has the RLP (right lifting property) with respect to $i$.
\end{defn}

We will use the (closed) Quillen model category axiom formulation of \cite{Hirschhorn}, except that we do not require functorial factorization of maps in the axiom M5. 

\begin{defn}[Quillen model categories]
\label{defn:modelcat}
A Quillen model category (${\cat M}$, ${\cat W}$, ${\cat Cof}$, ${\cat Fib}$) consists of a category ${\cat M}$ and three distinguished classes of maps ${\cat W}$, ${\cat Cof}$, ${\cat Fib}$ - the weak equivalences, the cofibrations and respectively the fibrations of ${\cat M}$, subject to the axioms below:
\begin{description}
\item[M1] 
${\cat M}$ is complete and cocomplete. 
\item[M2] 
(Two out of three axiom) If $f, g$ are maps of ${\cat M}$ such that $gf$ is defined, and if two of $f, g, gf$ are weak equivalences, then so is the third.
\item[M3] 
(Retract axiom) Weak equivalences, cofibrations and fibrations are closed under retracts.
\item[M4]
(Lifting axiom) Cofibrations have the LLP with respect to trivial fibrations, and trivial cofibrations have the LLP with respect to fibrations.
\item[M5]
(Factorization axiom) Any map of ${\cat M}$ admits a factorization as a cofibration followed by a trivial fibration, and a factorization as a trivial cofibration followed by a fibration.
\end{description}
\end{defn}

The axiom M3 states that given a commutative diagram
\begin{center}
$\xymatrix{
	A \ar[r] \ar[d]_f \ar@/^1pc/[rr]^{1_A} &
	C \ar[d]^g \ar[r] & A \ar[d]^f \\
	B \ar[r] \ar@/_1pc/[rr]_{1_B} &
	D \ar[r] & B
	}$
\end{center}
if $g$ is a weak equivalence (resp. cofibration, resp. fibration) then so is its retract $f$.

\begin{prop}
\label{prop:llpcharact}
In a Quillen model category any two of the following classes of maps of ${\cat M}$ - the cofibrations, the trivial cofibrations, the fibrations and the trivial fibrations - determine each other by the following rules: a map is
\begin{enumerate}
\item
A cofibration $\Leftrightarrow$ it has the LLP with respect to all trivial fibrations
\item
A trivial cofibration $\Leftrightarrow$ it has the LLP with respect to all fibrations
\item
A fibration $\Leftrightarrow$ it has the RLP with respect to all trivial cofibrations
\item
A trivial fibration $\Leftrightarrow$ it has the RLP with respect to all cofibrations $\square$
\end{enumerate}
\end{prop}

\begin{proof}
This is a direct consequence of the axioms M1-M5.
\end{proof}

\begin{prop}
\label{prop:relclosedmodcat}
Any Quillen model category is an ABC model category. 
\end{prop}

\begin{proof} A Quillen model category trivially satisfies the axioms CF2 and CF4. The axioms CF1, CF3, CF5 and CF6 are satisfied as a consequence of \propositionref{prop:llpcharact}. A dual argument shows that a Quillen model category satisfies the axioms F1-F6.
\end{proof}

If ${\cat M}$ is a Quillen model category and ${\cat D}$ is a small category, then the category of diagrams ${\cat M}^{\cat D}$ does not generally form a Quillen model category (except in important particular cases, for example when ${\cat M}$ is cofibrantly generated or when ${\cat D}$ is a Reedy category). But we will see further down (\theoremref{thm:generalpointwisecofstructure}) that ${\cat M}^{\cat D}$ carries an ABC model category structure, and in that sense one can always 'do homotopy theory' on ${\cat M}^{\cat D}$.

A Quillen model category ${\cat M}$ is called {\it left proper} if for any pushout diagram
\begin{center}
$ \xymatrix {
	A \ar[r]^j \ar@{>->}[d]_i &
	C \ar@{>-->}[d] \\
	B \ar@{-->}[r]^{j^{'}} &
	D
	}$
\end{center}
with $i$ a cofibration and $j$ a weak equivalence, the map $j^{'}$ is a weak equivalence. ${\cat M}$ is called {\it right proper} if for any pullback diagram
\begin{center}
$ \xymatrix {
	D \ar@{-->}[r]^{q^{'}} \ar@{-->>}[d] &
	B \ar@{->>}[d]^p \\
	C \ar[r]^q &
	A
	}$
\end{center}
with $p$ a fibration and $q$ a weak equivalence, the map $q^{'}$ is a weak equivalence. ${\cat M}$ is called {\it proper} if it is left and right proper.

From this definition and from \propositionref{prop:relclosedmodcat} we immediately get
\begin{prop}
\label{prop:relproperclosedmodcat}
Any proper Quillen model category is a proper ABC model category. $\square$
\end{prop}

\section{Baues cofibration categories}
%\index{Baues (co)fibration category}
We next turn our attention to the notion of (co)fibration category as defined by Baues. We state below the axioms of a Baues cofibration category, in a slightly modified but equivalent form to \cite{Baues1}, Sec. 1.1.

\begin{defn}
\label{defn:bauescofcat}
A Baues cofibration category $({\cat M}$, ${\cat W}$, ${\cat Cof})$ consists of a category ${\cat M}$ and two distinguished classes of maps ${\cat W}$ and ${\cat Cof}$ - the weak equivalences and the cofibrations of ${\cat M}$ - subject to the axioms below:
\begin{description}
\item[BCF1] 
All isomophisms of ${\cat M}$ are trivial cofibrations. Cofibrations are stable under composition. 
\item[BCF2] 
(Two out of three axiom) If $f, g$ are maps of ${\cat M}$ such that $gf$ is defined, and if two of $f, g, gf$ are weak equivalences, then so is the third.
\item[BCF3] 
(Pushout and excision axiom) Given a solid diagram in ${\cat M}$, with $i$ a cofibration,
\begin{center}
$ \xymatrix {
	A \ar[r]^f \ar@{>->}[d]_i &
	C \ar@{>-->}[d]^j \\
	B \ar@{-->}[r]^g &
	D
	}$
\end{center}
then the pushout exists in ${\cat M}$ and $j$ is a cofibration. Moreover:
\begin{enumerate}
\item[(a)]
If $i$ is a trivial cofibration, then $j$ is a trivial cofibration
\item[(b)]
If $f$ is a weak equivalence, then $g$ is a weak equivalence.
\end{enumerate}
\item[BCF4]
(Factorization axiom) Any map of ${\cat M}$ admits a factorization as a cofibration followed by a weak equivalence.
\item[BCF6]
(Axiom on fibrant models) For each object $A$ of ${\cat M}$ there is a trivial cofibration $A \ra B$, with $B$ satisfying the property that each trivial cofibration $C \ra B$ admits a left inverse.
\end{description}
\end{defn}

A Baues cofibration category may not have  an initial object, but if it does then by BCF1 the initial object is cofibrant. We will not state the axioms for a Baues fibration category - they are dual to the above axioms.

\begin{prop}[Relation with Baues cofibration categories]
\label{prop:relcatcofcatbaues}
\mbox{}
\begin{enumerate}
\item
Any Baues cofibration category with an initial object is a left proper precofibration category.
\item
Any Baues fibration category with a terminal object is a right proper prefibration category.
\end{enumerate}
\end{prop}

\begin{proof}
Easy consequence of the axioms and of \lemmaref{lem:equivaxiomcf3}.
\end{proof}

\section{Waldhausen categories}
%\index{Waldhausen (co)fibration category}
We list below the axioms we'll use for a Waldhausen cofibration category. These axioms are equivalent to the axioms Cof1-Cof3, Weq1 and Weq2 of \cite{Waldhausen1}. Axioms for a Waldhausen fibration category will of couse be dual to the axioms below. 
\begin{defn}
\label{defn:waldcofcat}
A Waldhausen cofibration category $({\cat M}$, ${\cat W}$, ${\cat Cof})$ consists of a category ${\cat M}$ and two distinguished classes of maps ${\cat W}$ and ${\cat Cof}$ - the weak equivalences and the cofibrations of ${\cat M}$, subject to the axioms below:
\begin{description}
\item[WCF1] 
${\cat M}$ is pointed. All isomophisms of ${\cat M}$ are trivial cofibrations. All objects of ${\cat M}$ are cofibrant. Cofibrations are stable under composition. 
\item[WCF2] 
(Pushout axiom) Given a solid diagram in ${\cat M}$, with $i$ a cofibration,
\begin{center}
$ \xymatrix {
	A \ar[r] \ar@{>->}[d]_i &
	C \ar@{>-->}[d]^j \\
	B \ar@{-->}[r] &
	D
	}$
\end{center}
then the pushout exists in ${\cat M}$ and $j$ is a cofibration
\item[WCF3]
(Gluing axiom) In the diagram below
\begin{center}
$ \xymatrix {
        A_1 \ar@{>->}[rr]^{f_{12}} \ar[dd]_{u_1} \ar[rd]_{f_{13}} &  & A_2 \ar'[d][dd]^{u_2} \ar[rd] & \\
        & A_3 \ar@{>->}[rr] \ar[dd]^(.3){u_3} & & A_4 \ar[dd]^{u_4} \\
        B_1 \ar@{>->}'[r][rr]_{g_{12}} \ar[rd]_{g_{13}} & & B_2 \ar[rd] & \\
        & B_3 \ar@{>->}[rr] & & B_4
	}$
\end{center}
if $f_{12}, g_{12}$ are cofibrations, $u_1, u_2, u_3$ are weak equivalences and the top and bottom faces are pushouts, then $u_4$ is a weak equivalence.
\end{description}
\end{defn}

We have the following

\begin{prop}[Relation with Waldhausen cofibration categories]
\label{prop:relcatcofcatwaldhausen}
\mbox{}
\begin{enumerate}
\item
\begin{enumerate}
\item
If ${\cat M}$ is a pointed precofibration category, then ${\cat M}_{cof}$ is a Waldhausen cofibration category.
\item
If ${\cat M}$ is a Waldhausen cofibration category satisfying the 2 out of 3 axiom CF2 and the cylinder axiom CFObj4, then it is a precofibration category.
\end{enumerate}
\item
\begin{enumerate}
\item
If ${\cat M}$ is a pointed prefibration category, then ${\cat M}_{fib}$ is a Waldhausen fibration category.
\item
If ${\cat M}$ is a Waldhausen fibration category satisfying the 2 out of 3 axiom F2 and the path object axiom FObj4, then it is an prefibration category.
\end{enumerate}
\end{enumerate}
\end{prop}

\begin{proof}
Part (1) (a) follows from the Gluing Lemma \ref{lem:gluinglemma1}, part (1) (b) from \lemmaref{lem:equivaxiomcf3} and \propositionref{prop:relcatcofobj}, and part (2) is dual to the above.
\end{proof}

Waldhausen categories have been further studied by Thomason-Trobaugh \cite{Thomason-Trobaugh}, Weiss and Williams \cite{Weiss1}, \cite{Weiss2}, \cite{Weiss3}.